\part{Глава №1. Математические основы функционирования ЭВМ.}
\chapter{Тема №1. Системы счисления.}
\section{Лекция №1. Позиционные и непозиционные сс. Перевод из 10–сс в N–сс и обратно. Быстрый перевод из 2–сс в 4–, 8– и 16–сс. Базовая целочисленная арифметика в N–сс.}

\subsection{Что такое система счисления}
\textbf{Система счисления (СС)} — совокупность приёмов и правил для записи чисел с помощью набора символов (цифр).

Различают \textbf{непозиционные} и \textbf{позиционные} системы.

\subsubsection*{Непозиционные системы счисления}
Значение цифры не зависит от её места в записи числа.
Примеры:
\begin{itemize}
    \item \textbf{Унарная} (единичная): число $N$ обозначается $N$ одинаковыми знаками (например, палочки: $||| = 3$).
    \item \textbf{Римская}: $I=1$, $V=5$, $X=10$, $L=50$, $C=100$, $D=500$, $M=1000$. Правила сложения/вычитания, но позиция влияет не так жёстко, как в позиционных, — нет веса разряда, кратного основанию.
\end{itemize}
Недостатки: громоздкость, трудность арифметических операций. Сейчас почти не используются в вычислениях.

\subsubsection*{Позиционные системы счисления}
Значение цифры зависит от её \textbf{позиции} (разряда).

Любое число $A$ в позиционной системе с основанием $p$ записывается как последовательность цифр:
\[
A = a_{n-1} a_{n-2} \dots a_1 a_0 \, . \, a_{-1} a_{-2} \dots a_{-m}
\]
Его числовое значение определяется по формуле:
\[
A = \sum_{i=-m}^{n-1} a_i \cdot p^i
\]
Где:
\begin{itemize}
    \item $p$ — \textbf{основание системы} (целое, $p \ge 2$),
    \item $a_i$ — цифры ($0 \le a_i \le p-1$),
    \item $n$ — количество разрядов целой части,
    \item $m$ — количество разрядов дробной части.
\end{itemize}

\textbf{Примеры позиционных систем:}
\begin{itemize}
    \item Десятичная ($p = 10$, алфавит 0--9)
    \item Двоичная ($p = 2$, алфавит 0, 1)
    \item Троичная ($p = 3$, алфавит 0, 1, 2)
    \item Восьмеричная ($p = 8$, алфавит 0--7)
    \item Шестнадцатеричная ($p = 16$, алфавит 0--9, A=10, B=11, C=12, D=13, E=14, F=15)
\end{itemize}

В нашем курсе особое значение имеют двоичная, восьмеричная и шестнадцатеричная системы из-за их прямой связи с организацией памяти.

\subsection{Перевод чисел из десятичной системы в систему с основанием N}

\subsubsection*{Перевод целой части}
\textbf{Алгоритм:} последовательно делить число на основание $N$, запоминая остатки. Остатки, выписанные в обратном порядке, дают искомую запись.

\textit{Пример:} переведём $156_{10}$ в двоичную систему ($N=2$).
\begin{verbatim}
156 / 2 = 78 (остаток 0)
78  / 2 = 39 (остаток 0)
39  / 2 = 19 (остаток 1)
19  / 2 = 9  (остаток 1)
9   / 2 = 4  (остаток 1)
4   / 2 = 2  (остаток 0)
2   / 2 = 1  (остаток 0)
1   / 2 = 0  (остаток 1)
\end{verbatim}
Результат: $10011100_2$.

\textit{Пример:} $156_{10} \to 16$-я система.
\begin{verbatim}
156 / 16 = 9 (остаток 12 -> C)
9   / 16 = 0 (остаток 9)
\end{verbatim}
Результат: $9C_{16}$.

\subsubsection*{Перевод дробной части (упоминание)}
Хотя в первой теме акцент на целых, полезно знать: дробную часть умножают на $N$, записывают целую часть результата, затем отбрасывают её и повторяют с оставшейся дробной частью. Конечность или периодичность зависит от числа и основания.

\textit{Пример:} $0.375_{10} \to 2$-я система.
\[
0.375 \times 2 = 0.75 \text{ (целая 0)},\quad 0.75 \times 2 = 1.5 \text{ (целая 1)},\quad 0.5 \times 2 = 1.0 \text{ (целая 1)}
\]
Получаем $0.011_2$.

\subsection{Перевод из системы с основанием N в десятичную}
Используем развёрнутую запись по степеням основания.

\textit{Пример:} $10011100_2$ в десятичную.
\[
10011100_2 = 1\cdot2^7 + 0\cdot2^6 + 0\cdot2^5 + 1\cdot2^4 + 1\cdot2^3 + 1\cdot2^2 + 0\cdot2^1 + 0\cdot2^0
= 128 + 16 + 8 + 4 = 156_{10}.
\]

\textit{Пример:} $9C_{16} = 9\cdot16^1 + 12\cdot16^0 = 144 + 12 = 156_{10}$.

\subsection{Быстрый перевод между двоичной, восьмеричной и шестнадцатеричной системами}
Поскольку $8 = 2^3$, а $16 = 2^4$, перевод сводится к группировке двоичных разрядов:
\begin{itemize}
    \item \textbf{Из двоичной в восьмеричную:} разбить двоичную запись справа налево на группы по 3 бита (\textbf{триады}), каждую группу заменить восьмеричной цифрой.
    \item \textbf{Из восьмеричной в двоичную:} каждую цифру заменить тремя двоичными разрядами.
    \item \textbf{Из двоичной в шестнадцатеричную:} разбить на группы по 4 бита (\textbf{тетрады}), заменить 16-ричной цифрой.
    \item \textbf{Из 16-ричной в двоичную:} каждую цифру заменить четырьмя битами.
\end{itemize}

\textit{Пример:} $110111001010_2$.
Для 8-й: \texttt{110 111 001 010} $\to$ \texttt{6 7 1 2} $\to$ $6712_8$.
Для 16-й: \texttt{1101 1100 1010} $\to$ \texttt{D C A} $\to$ $DCA_{16}$.
Обратно: $DCA_{16} \to D=1101, C=1100, A=1010 \to 110111001010_2$.

Если в старшей группе не хватает битов, дополняем слева нулями.

Перевод между 8-й и 16-й проще всего делать через двоичную как промежуточную.

\subsection{Базовая целочисленная арифметика в N-й системе счисления}
Арифметические действия в любой позиционной системе выполняются по тем же правилам, что и в десятичной, но с учётом \textbf{основания}.

\subsubsection*{Сложение}
Складываем поразрядно справа налево. Если сумма в разряде $\ge p$, вычитаем $p$ и делаем перенос 1 в старший разряд.

\textit{Пример в двоичной ($p=2$):}
\begin{verbatim}
  1011_2  (11)
+ 1101_2  (13)
---------
=11000_2  (24)
\end{verbatim}
Столбиком: $1+1=10_2$ (0 пишем, 1 перенос); следующий: $1+0+1=10_2$ (0, 1 перенос); следующий: $0+1+1=10_2$ (0, 1 перенос); старший: $1+1+1=11_2$ (1, 1 перенос). Итог: $11000_2$.

\subsubsection*{Вычитание}
Занимаем единицу из старшего разряда, которая превращается в $p$ в текущем разряде. Далее вычитаем с учётом занятого.

\textit{Пример в 16-й:} $A3_{16} - 4F_{16}$.
$A=10$, $3=3$. Младший: $3 - F(15)$ нельзя, занимаем у $A$: остаётся 9, в младший приходит $16+3=19$. $19-15=4$. Старший: $9-4=5$. Итог: $54_{16}$ ($84_{10}$).

\subsubsection*{Умножение}
Используем таблицу умножения для данного основания. Умножаем каждый разряд, результаты сдвигаем и складываем в данной СС.

\textit{Пример в 8-й:} $23_8 \times 5_8$.
$3\times5=15_{10}=17_8$ (7 пишем, 1 в уме). $2\times5=10_{10}+1=11_{10}=13_8$. Итого $137_8$. Проверка: $23_8=19_{10}$, $5_8=5$, $19\times5=95_{10}=137_8$.

\subsubsection*{Деление}
Деление уголком по тем же принципам, но подбор цифр и вычитание — в данной СС.




\newpage
\section{Семинар №1. Перевод из 10–сс в N–сс и обратно. Быстрый перевод из 2–сс в 4–, 8– и 16–сс. Базовая целочисленная арифметика в N–сс.}

\subsection*{Задача 1 (перевод целых чисел $10 \to N$)}
Переведите число $250_{10}$ в двоичную, восьмеричную и шестнадцатеричную системы.

\textbf{Решение:}
\begin{itemize}
    \item В двоичную: $250 / 2 = 125 (0)$, $125/2=62 (1)$, $62/2=31 (0)$, $31/2=15 (1)$, $15/2=7 (1)$, $7/2=3 (1)$, $3/2=1 (1)$, $1/2=0 (1)$. Собираем снизу: $11111010_2$.
    \item В восьмеричную: $250 / 8 = 31\ (\text{ост. }2)$, $31/8=3\ (\text{ост. }7)$, $3/8=0\ (\text{ост. }3)$ $\to 372_8$.
    \item В шестнадцатеричную: $250 / 16 = 15\ (\text{ост. }10=A)$, $15/16=0\ (\text{ост. }15=F)$ $\to FA_{16}$.
\end{itemize}

\subsection*{Задача 2 (перевод $N \to 10$)}
Даны числа: $1100101_2$, $145_8$, $A9_{16}$. Переведите их в десятичную систему.

\textbf{Решение:}
\begin{align*}
1100101_2 &= 1\cdot64 + 1\cdot32 + 0\cdot16 + 0\cdot8 + 1\cdot4 + 0\cdot2 + 1\cdot1 = 64+32+4+1 = 101_{10}.\\
145_8     &= 1\cdot64 + 4\cdot8 + 5\cdot1 = 64+32+5 = 101_{10}.\\
A9_{16}   &= 10\cdot16 + 9 = 160+9 = 169_{10}.
\end{align*}

\subsection*{Задача 3 (быстрый перевод $2\leftrightarrow 8$, $2\leftrightarrow 16$)}
Дано двоичное число $110111001010_2$. Переведите его в восьмеричную и шестнадцатеричную системы.

\textbf{Решение:}
Группы по 3 (8-я): \texttt{110 111 001 010} $\to$ \texttt{6 7 1 2} $\to 6712_8$.\\
Группы по 4 (16-я): \texttt{1101 1100 1010} $\to$ \texttt{D C A} $\to DCA_{16}$.

Дано $2F5_{16}$. Переведите в двоичную и восьмеричную.

\textbf{Решение:}
$2=0010$, $F=1111$, $5=0101$ $\to 001011110101_2 \to$ убираем ведущие нули: $1011110101_2$.\\
В 8-ю: разбиваем $1011110101_2$ на триады справа: \texttt{1 011 110 101} $\to$ \texttt{001 011 110 101} $\to$ \texttt{1 3 6 5} $\to 1365_8$.

\subsection*{Задача 4 (арифметика)}
Выполните действия в указанных системах. Проверьте переводом в десятичную.
\begin{enumerate}
    \item[a)] $1011_2 + 1101_2$
    \item[b)] $11001_2 - 1011_2$
    \item[c)] $23_8 \times 5_8$
    \item[d)] $1A3_{16} + 4F_{16}$
    \item[e)] $200_8 / 7_8$ (целочисленное деление)
\end{enumerate}

\textbf{Решение:}
\begin{enumerate}
    \item[a)] $1011_2 (11) + 1101_2 (13) = 11000_2 (24).$
    \item[b)] $11001_2 (25) - 1011_2 (11) = 1110_2 (14).$\\
    Процесс: занимаем. $11001 - 1011$ столбиком даёт $1110$.
    \item[c)] $23_8 \times 5_8 = 137_8$. Проверка: $23_8=19$, $5_8=5$, $19\times5=95$, $95_{10}=137_8$.
    \item[d)] $1A3_{16} = 1\cdot256 + 10\cdot16 + 3 = 419$. $4F_{16}=4\cdot16+15=79$. $419+79=498$. $498_{10}$ в 16-ю: $498/16=31\ (\text{ост. }2)$, $31/16=1\ (\text{ост. }15=F)$, $1/16=0\ (\text{ост. }1)$ $\to 1F2_{16}$. Итог: $1A3_{16}+4F_{16}=1F2_{16}$.
    \item[e)] $200_8 = 2\times64 + 0 = 128_{10}$. $7_8=7$. $128 / 7 = 18$ (ост. 2). $18_{10}$ в 8-ю: $22_8$. Остаток 2. Частное $22_8$, остаток 2.
\end{enumerate}


\newpage
\chapter{Тема №2. Представление чисел в памяти компьютера.}
\section{Лекция №2. Целые числа: знаковые и беззнаковые. Диапазоны. Прямой, обратный, сдвиговый, дополнительный коды. Оптимизация арифметики отрицательных чисел с помощью доп. кода. Побитовые операции. Двоичная арифметика.}

\subsection{Представление беззнаковых целых чисел (unsigned)}
В ячейках памяти информация хранится в битах. Для беззнаковых целых все биты представляют значение.

Если под число отведено $n$ бит, диапазон значений: $0 \dots 2^n-1$.

\textit{Пример:} $n=8$ (байт). Диапазон $0\dots255$. Число $153_{10} = 10011001_2$.

Формула: $\sum_{i=0}^{n-1} b_i \cdot 2^i$.

Арифметические операции над беззнаковыми: обычная двоичная арифметика, переполнение фиксируется флагом переноса \textbf{CF} (Carry Flag) при выходе за разрядную сетку.

\subsection{Представление знаковых целых чисел}
Для представления отрицательных чисел используют старший бит как \textbf{знаковый}: 0 — плюс, 1 — минус.
Остальные биты могут интерпретироваться по-разному. Существуют четыре основных кода: прямой, обратный, сдвиговый и дополнительный.

\subsubsection{Прямой код (Sign-Magnitude)}
Старший бит — знак, остальные — модуль числа.
\textit{Пример (8 бит):} $+18 = 00010010_2$, $-18 = 10010010_2$.

Диапазон: $-(2^{n-1}-1) \dots +(2^{n-1}-1)$. Два нуля: $+0$ ($00\dots0$) и $-0$ ($10\dots0$).
Арифметика требует анализа знаков, что неудобно для АЛУ.

\subsubsection{Обратный код (Ones' Complement)}
Отрицательное число получается побитовой инверсией модуля (все 0 заменяются 1, 1 на 0), включая знаковый бит.
\textit{Пример:} $+18 = 00010010_2$; $-18 = 11101101_2$.

Диапазон такой же, как у прямого. Два нуля: $+0 = 00\dots0$, $-0 = 11\dots1$.
Сложение проще, но возникает циклический перенос (end-around carry), что усложняет аппаратную реализацию.

\subsubsection{Сдвиговый код (Biased Representation, код со смещением)}
В сдвиговом коде значение числа $X$ представляется как беззнаковое целое $X + \text{bias}$, где bias — фиксированная константа (смещение). Для $n$-разрядной сетки обычно выбирают $\text{bias} = 2^{n-1}$ или $2^{n-1}-1$.

При $\text{bias}=2^{n-1}$ (например, $128$ для $n=8$) получается диапазон $X \in [-128, 127]$, которому соответствуют коды $00000000_2$ ($-128$) … $11111111_2$ ($+127$).
\textit{Пример ($n=8$, bias $=128$):} $+18 \to 18+128=146 = 10010010_2$; $-18 \to -18+128=110 = 01101110_2$.

Заметим, что нулевое значение представляется единственным образом: $0 \to 128 = 10000000_2$.
Сдвиговый код удобен для сравнения чисел (сравниваются как беззнаковые), но арифметические операции требуют коррекции результата на величину bias, поэтому в чистом виде для целочисленной арифметики почти не применяется. Широко используется для хранения порядка (экспоненты) чисел с плавающей точкой (стандарт IEEE 754).

\subsubsection{Дополнительный код (Two's Complement)}
Отрицательное число формируется путём инверсии модуля и прибавления 1 к младшему разряду (либо вычитанием модуля из $2^n$).
\textit{Пример ($n=8$):} $+18 = 00010010_2$. Инверсия: $11101101_2$, $+1 \to 11101110_2 = -18$.

\textbf{Альтернативно:} $2^8 - 18 = 256 - 18 = 238 = 11101110_2$.

Диапазон: $-2^{n-1} \dots 2^{n-1}-1$. Ноль единственный: $00\dots0$ (инверсия $11\dots1$, $+1 = 00\dots0$ с отбрасыванием переполнения).

\subsubsection{Оптимизация арифметики отрицательных чисел с помощью дополнительного кода}
Дополнительный код стал промышленным стандартом, потому что позволяет \textbf{выполнять вычитание как сложение} без анализа знаков и без циклического переноса:
$A - B = A + (-B)$.
Единый алгоритм работает и для знаковых, и для беззнаковых чисел. АЛУ строит только сумматор.

\textit{Пример:} $5 - 3 = 5 + (-3)$. $5=00000101_2$, $-3=11111101_2$ (доп. код). Сумма: $00000010_2 = 2$ (перенос за пределы разрядной сетки игнорируется).

При расширении разрядности (например, из 8 в 16 бит) используется \textbf{знаковое расширение} (sign extension): все новые старшие биты заполняются копией знакового бита.

\subsection{Побитовые операции}
\begin{itemize}
    \item \textbf{NOT (побитовое отрицание, инверсия)} $\sim$: $0\leftrightarrow 1$ для каждого бита.
    \item \textbf{AND (И)} $\&$: 1, если оба бита 1.
    \item \textbf{OR (ИЛИ)} $|$: 1, если хотя бы один бит 1.
    \item \textbf{XOR (исключающее ИЛИ)} $\wedge$: 1, если биты различны.
    \item \textbf{Сдвиги:}
    \begin{itemize}
        \item \textit{Логический сдвиг влево} (<<): биты сдвигаются влево, справа заполняются нулями. Умножение на $2^k$ (для беззнаковых).
        \item \textit{Логический сдвиг вправо} (>> для unsigned): сдвиг вправо, слева нули. Деление на $2^k$ (беззнаковое).
        \item \textit{Арифметический сдвиг вправо} (>> для signed): сдвиг вправо, слева копируется знаковый бит. Сохраняет знак при делении степеней двойки.
        \item \textit{Циклический сдвиг} (rotate): биты перемещаются, вышедшие за край вставляются с противоположной стороны.
    \end{itemize}
\end{itemize}

\textit{Примеры (8 бит):}
\begin{verbatim}
x = 0b01101011
x & 0b00001111 = 0b00001011  (выделение младших 4 бит)
x | 0b10000000 = 0b11101011  (установка старшего бита)
x ^ x = 0                     (обнуление)
x << 2 = 0b10101100           (умножение на 4, может вызвать переполнение)
\end{verbatim}

\subsection{Двоичная арифметика с учётом ограниченной разрядности}

\textbf{Сложение:} побитово с переносом.

\textbf{Переполнение (overflow)} при знаковых операциях фиксируется флагом \textbf{OF}. Пример: $127 + 1 = -128$ в 8-битном дополнительном коде ($01111111 + 00000001 = 10000000$).

\textbf{Перенос (carry)} для беззнаковых фиксируется флагом \textbf{CF}. $255+1=0$ с CF=1.

\textbf{Вычитание:} через дополнительный код. $A - B = A + (\sim B + 1)$.

\textbf{Умножение:} столбиком, в процессорах реализовано сдвигами и сложением. Умножение двух $n$-битных даёт $2n$ бит.

\textbf{Деление:} аппаратно или программно. Выдаёт частное и остаток.

\textbf{Флаги состояния (для понимания):}
\begin{itemize}
    \item ZF (Zero): результат = 0.
    \item SF (Sign): копия старшего бита результата (знак в доп. коде).
    \item CF (Carry): беззнаковое переполнение.
    \item OF (Overflow): знаковое переполнение.
\end{itemize}


\newpage
\section{Семинар №2. Целые числа: знаковые и беззнаковые. Диапазоны. Прямой, обратный, сдвиговый, дополнительный коды. Оптимизация арифметики отрицательных чисел с помощью доп. кода. Побитовые операции. Двоичная арифметика.}

\subsection*{Задача 1 (диапазоны)}
Сколько различных чисел можно представить 12-битным беззнаковым целым? А 12-битным знаковым в дополнительном коде? Укажите минимальное и максимальное значение для каждого случая.

\textbf{Решение:}
\begin{itemize}
    \item Беззнаковое 12 бит: $2^{12} = 4096$ чисел; диапазон $[0, 4095]$.
    \item Знаковое доп. код: $2^{12} = 4096$; диапазон $[-2048, 2047]$ (от $-2^{11}$ до $2^{11}-1$).
\end{itemize}

\subsection*{Задача 2 (прямой, обратный, сдвиговый, дополнительный коды)}
Представьте число $-37$ в 8-битных прямом, обратном, сдвиговом (со смещением 128) и дополнительном кодах.

\textbf{Решение:}
Модуль $37_{10} = 00100101_2$ (дополняем до 8 бит).
\begin{itemize}
    \item Прямой код: знак 1, модуль: $10100101_2$.
    \item Обратный: инвертируем все биты модуля вместе со знаковым (получаем инверсию $+37$): $+37 = 00100101_2 \to \text{инверсия } 11011010_2$. Это и есть $-37$ в обратном коде.
    \item Сдвиговый (bias = 128): значение $X = -37$, код = $X + 128 = 91$. $91_{10} = 01011011_2$ (дополняем до 8 бит). Ответ: $01011011_2$.
    \item Дополнительный: обратный код $+1$: $11011010_2 + 1 = 11011011_2$.
\end{itemize}

Проверка дополнительного: $256 - 37 = 219 = 11011011_2$.

\subsection*{Задача 3 (сложение в дополнительном коде)}
Выполните сложение 8-битных чисел в дополнительном коде, интерпретируйте результат:
\begin{enumerate}
    \item[a)] $45 + (-23)$
    \item[b)] $100 + 30$ (проверьте флаги CF и OF)
    \item[c)] $(-128) + (-1)$
\end{enumerate}

\textbf{Решение:}
\begin{enumerate}
    \item[a)] $45 = 00101101_2$; $-23$: $23=00010111_2$, инв.$+1=11101001_2$. Сумма: $00101101+11101001=1\underline{00010110}$ (9 бит). Отбрасываем 9-й бит $\to 00010110_2 = 22_{10}$. Верно.
    \item[b)] $100 = 01100100_2$, $30 = 00011110_2$. Сумма $= 10000010_2$ ($130_{10}$ беззнак., знаковое: $-126$). CF=0 (нет переноса из старшего бита), OF=1 (перенос в знаковый разряд был 1, перенос из знакового — 0).
    \item[c)] $-128 = 10000000_2$, $-1 = 11111111_2$. Сумма $= 1\,01111111$ (9 бит). 8-й бит отбрасываем $\to 01111111_2 = 127$. OF=1 (два отрицательных дали положительное), CF=1.
\end{enumerate}

\subsection*{Задача 4 (побитовые операции)}
Дано $x = 0x5A$ (hex), $y = 0x3C$. Вычислите (в двоичном и hex виде):
\begin{enumerate}
    \item[a)] $x$ \& $y$
    \item[b)] $x$ | $y$
    \item[c)] $x$ $\wedge$ $y$
    \item[d)] $\sim x$ (8-битное)
    \item[e)] $x << 2$ (логический сдвиг)
    \item[f)] арифметический сдвиг вправо для знакового $x$ (8 бит, $x = 0xA6$) на 2.
\end{enumerate}

\textbf{Решение:}
$x = 0x5A = 01011010_2$, $y = 0x3C = 00111100_2$.
\begin{enumerate}
    \item[a)] AND: $01011010 \,\&\, 00111100 = 00011000_2 = 0x18$.
    \item[b)] OR: $01011010 \,|\, 00111100 = 01111110_2 = 0x7E$.
    \item[c)] XOR: $01011010 \,\wedge\, 00111100 = 01100110_2 = 0x66$.
    \item[d)] $\sim x$ (8 бит) = $10100101_2 = 0xA5$.
    \item[e)] $x << 2$: $01011010 << 2 = 01101000_2 = 0x68$ (вытеснение старшего бита, результат $01101000$).
    \item[f)] $x = 0xA6 = 10100110_2$ (знаковый бит 1). Арифм. сдвиг вправо на 2: $10100110 \gg 1 = 11010011$; $\gg 1 = 11101001$. Итог: $11101001_2 = 0xE9$.
\end{enumerate}

\subsection*{Задача 5 (флаги и интерпретации)}
8-битные числа (дополнительный код). Выполните вычитание $80 - 50$ двумя способами: непосредственно в столбик (если умеете) и через сложение с доп. кодом. Укажите флаги результата.

\textbf{Решение:}
$80 = 01010000_2$, $50 = 00110010_2$.
Вычитание: $A - B = A + (-B)$. $-50 =$ инв$(00110010)+1 = 11001101+1=11001110_2$.
Сумма: $01010000 + 11001110 = 1\,00011110$ (9 бит). Отбросив перенос, получаем $00011110_2 = 30$.
Флаги (в 8-битном контексте): SF=0, ZF=0, CF=1 (беззнаковое $80+206=286>255$, перенос), OF=0 (перенос в знаковый разряд и из него совпадают: оба 1).

\subsection*{Задача 6 (арифметика в сдвиговом коде)}
Два числа $A=-20$ и $B=35$ представлены в 8-битном сдвиговом коде со смещением 128. Выполните их сложение в этом представлении и объясните, как получить истинную сумму в дополнительном коде.

\textbf{Решение:}
Сначала представим числа в сдвиговом коде с bias = 128:
\begin{itemize}
    \item $A=-20 \to \text{code}_A = -20 + 128 = 108 = 01101100_2$.
    \item $B=35 \to \text{code}_B = 35 + 128 = 163 = 10100011_2$.
\end{itemize}

Сложение кодов как беззнаковых чисел:
$01101100_2 + 10100011_2 = 1\,00001111_2$ (9 бит). Отбрасываем старший (девятый) бит, получаем $00001111_2 = 15_{10}$.

Этот код соответствует значению $15 - 128 = -113$ в сдвиговом коде, но мы ожидаем сумму $A+B = 15$. Ошибка возникла потому, что при сложении двух смещённых чисел получается $(A+128)+(B+128) = A+B+256$, что в 8-битной сетке даёт $A+B+256 \pmod{256} = A+B$ (так как $256 \equiv 0$). Действительно, $108+163 = 271$, $271 \bmod 256 = 15$, а 15 в сдвиговом коде соответствует $15 - 128 = -113$, что не равно 15. Однако если мы интерпретируем результат как \textbf{беззнаковое} число 15 и затем вычтем bias, получим $-113$ — неверно.

Правильный способ сложения в сдвиговом коде: перед сложением вычесть bias из одного из операндов, либо после сложения вычесть bias. Формула:
\[
\text{code}_{A+B} = \text{code}_A + \text{code}_B - \text{bias}
\]
У нас $108 + 163 - 128 = 143$. $143$ в 8 битах $= 10001111_2$, что соответствует значению $143-128=15$ — верная сумма.

Таким образом, арифметика в сдвиговом коде сложнее, что объясняет, почему для целочисленных вычислений предпочитают дополнительный код.


\newpage
\section{Лекция №3. Числа с плавающей точкой. Мантисса и экспонента. Нормальная и нормализованная формы. Стандарт IEEE 754. Машинная точность.}

\subsubsection{Представление вещественных чисел: нормальная форма}
Вещественное число можно записать как:
\[
X = \pm M \times B^E
\]
где $M$ — \textbf{мантисса} (значащая часть), $B$ — \textbf{основание} (обычно 2 для компьютеров), $E$ — \textbf{экспонента} (порядок).

Чтобы запись была однозначной, мантиссу приводят к \textbf{нормализованной форме}: $1 \le |M| < B$. В двоичной системе это означает, что целая часть мантиссы равна 1 (кроме нуля и денормализованных чисел). Соответственно, в памяти хранятся:
\begin{itemize}
    \item Знак (1 бит)
    \item Экспонента со смещением
    \item Дробная часть мантиссы (без ведущей единицы, которая подразумевается)
\end{itemize}

\subsubsection{Стандарт IEEE 754}
Определяет форматы чисел с плавающей точкой, операции, округления, особые значения.

\paragraph{Одинарная точность (float, 32 бита)}
\[
\begin{array}{|c|c|c|}
\hline
\text{Знак (1 бит)} & \text{Экспонента (8 бит)} & \text{Мантисса (23 бита)} \\
\hline
\end{array}
\]
Порядок записан со \textbf{смещением} (bias) = 127. Истинная экспонента $E = e - \text{bias}$, где $e$ — хранимое беззнаковое целое.
\begin{itemize}
    \item Если $0 < e < 255$: \textbf{нормализованное число}. Значение $= (-1)^S \times (1.\text{дробная\_часть}) \times 2^{e-127}$.
    \item $e = 0$, мантисса $= 0$: $\pm 0$.
    \item $e = 0$, мантисса $\neq 0$: \textbf{денормализованные числа}. Значение $= (-1)^S \times (0.\text{дробная\_часть}) \times 2^{-126}$. Позволяют плавный переход к нулю (underflow).
    \item $e = 255$, мантисса $= 0$: $\pm\infty$ (переполнение, деление на 0).
    \item $e = 255$, мантисса $\neq 0$: NaN (Not a Number), например результат $\sqrt{-1}$.
\end{itemize}

\paragraph{Двойная точность (double, 64 бита)}
\[
\begin{array}{|c|c|c|}
\hline
\text{Знак (1)} & \text{Экспонента (11 бит)} & \text{Мантисса (52 бита)} \\
\hline
\end{array}
\]
Смещение bias = 1023. $e = 0\dots2046$ нормализованные, $e=0$ денормализованные, $e=2047$ особые случаи.

\subsubsection{Примеры кодирования}
Представим $-12.625$ в float IEEE 754.
\begin{enumerate}
    \item Переведём модуль в двоичную: $12_{10} = 1100_2$; $0.625 = 0.101_2$. Итого $1100.101_2$.
    \item Нормализуем: $1.100101_2 \times 2^3$.
    \item Экспонента истинная $= 3$, $e = 3 + 127 = 130 = 10000010_2$.
    \item Мантисса: отбрасываем $1.$ $\to$ $10010100000000000000000$ (23 бита).
    \item Знак $S = 1$.
\end{enumerate}
Собираем: \texttt{1 10000010 10010100000000000000000}.

\subsubsection{Машинная точность (машинный эпсилон)}
\textbf{Машинный эпсилон $\varepsilon$} — это разница между $1.0$ и наименьшим числом, б\'{о}льшим $1.0$, которое представимо в данном формате.

Для float: $1.0 = 1.000\ldots0 \times 2^0$. Следующее представимое число получается увеличением мантиссы на единицу младшего разряда, что соответствует значению $2^{-23} \approx 1.19\times10^{-7}$.

Для double: $\varepsilon = 2^{-52} \approx 2.22\times10^{-16}$.

$\varepsilon$ характеризует относительную плотность чисел: любое представимое число $x$ и следующее за ним отличаются примерно в $(1+\varepsilon)$ раз.

\subsection{Базовая арифметика чисел с плавающей точкой}

Арифметические операции с плавающей точкой сложнее целочисленных из-за необходимости выравнивания порядков, нормализации результата и округления. Рассмотрим каждый тип операций детально, вручную имитируя поведение float32 (23 бита дробной части).

\subsubsection{Сложение и вычитание}
Пусть даны два нормализованных числа: $X = M_x \times 2^{E_x}$, $Y = M_y \times 2^{E_y}$.

\paragraph*{Алгоритм сложения}
\begin{enumerate}
    \item \textbf{Выравнивание порядков}: меньший порядок приводится к большему путём сдвига мантиссы вправо с увеличением экспоненты. Если разница порядков $\ge 25$ (для float32, с учётом скрытого бита), мантисса меньшего числа при сдвиге станет нулевой — произойдёт \textbf{поглощение}.
    \item \textbf{Сложение мантисс} как целых чисел с учётом знаков. Мантиссы представляются с учётом скрытой единицы и знака.
    \item \textbf{Нормализация результата}: возможно, понадобится сдвиг мантиссы влево или вправо на 1 разряд с соответствующей коррекцией экспоненты.
    \item \textbf{Округление}: если точная мантисса имеет больше значащих разрядов, чем позволяет поле $f$, выполняется округление до ближайшего (по умолчанию к ближайшему чётному).
\end{enumerate}

\paragraph*{Пример сложения (точное)}

Сложим $A = 6.5_{10}$ и $B = 0.75_{10}$ в float32.

1. Представление:
\[
A = 6.5 = 110.1_2 = 1.101_2 \times 2^2 \quad (E_a = 2, M_a = +1.101_2)
\]
\[
B = 0.75 = 0.11_2 = 1.1_2 \times 2^{-1} \quad (E_b = -1, M_b = +1.1_2)
\]

2. Выравнивание: приводим к большему порядку $E = 2$. Мантиссу $B$ сдвигаем вправо на $2 - (-1) = 3$ разряда:
\[
M_b' = 0.0011_2 \quad (\text{сдвиг }1.1_2 \text{ на 3 вправо})
\]

3. Сложение мантисс:
\[
\begin{array}{r}
  1.1010_2 \quad (\text{дополним } M_a \text{ нулями до 5 бит})\\
+ 0.0011_2 \\
\hline
  1.1101_2 \\
\end{array}
\]
Сумма $= 1.1101_2 \times 2^2$.

4. Нормализация не требуется (мантисса в $[1,2)$). Результат $= 1.1101_2 \times 2^2 = 111.01_2 = 7.25_{10}$.

\paragraph*{Пример с округлением}

Пусть разрядность дробной части ограничена 2 битами (аналогично 23-м для float, но для простоты). Сложим $1.000_2 \times 2^0$ и $1.001_2 \times 2^{-3}$.
\[
A = 1.00_2 \times 2^0 \quad B = 1.001_2 \times 2^{-3} = 0.001001_2 \times 2^0
\]
После выравнивания ($B' = 0.0010_2$ при 4 битах дробной части) сумма $= 1.0010_2$. Если у нас мантисса хранит только 2 бита дробной части, то надо округлить $1.0010$ до ближайшего. В $1.0010$ третий бит после запятой = 1, а дальше нули — округляем до чётного: $1.00$ (так как ближайшие $1.00$ и $1.01$, $1.0010$ ровно посередине, а $1.00$ чётный). Результат $1.00 \times 2^0 = 1.0$.

\paragraph*{Пример вычитания с катастрофической потерей точности}

Вычтем два близких числа: $A = 1.001 \times 2^0$, $B = 1.000 \times 2^0$ (пусть 4 бита мантиссы). После выравнивания $B$ не меняется. Вычитаем мантиссы:
\[
1.001_2 - 1.000_2 = 0.001_2
\]
Теперь нужно нормализовать: $0.001_2 \times 2^0 = 1.000_2 \times 2^{-3}$. Видим, что точность мантиссы вроде сохранилась, но если исходные числа были приближёнными, относительная погрешность резко возрастает.

Более драматичный пример: $A = 1.23456789$, $B = 1.23456780$, разность $= 9 \times 10^{-8}$, но если числа хранятся с погрешностью $\sim 10^{-7}$, то разность может быть полностью потеряна.

\subsubsection{Умножение}
\[
X \times Y = (M_x \times M_y) \times 2^{E_x+E_y}
\]
\paragraph*{Алгоритм умножения}
\begin{enumerate}
    \item \textbf{Перемножение мантисс} как целых чисел (с учётом скрытых единиц). Если исходные мантиссы $1.f_x$ и $1.f_y$, то произведение $P = (1 + f_x) \times (1 + f_y) = 1 + f_x + f_y + f_x f_y$. Возможен дополнительный перенос в целую часть (результат может быть в диапазоне $[1,4)$).
    \item \textbf{Сложение экспонент}: $E_{\text{sum}} = E_x + E_y$ (истинных, а не хранимых со смещением; при реализации вычитают bias).
    \item \textbf{Нормализация}: если произведение мантисс $\ge 2$, то сдвигаем мантиссу вправо на 1 и увеличиваем экспоненту.
    \item \textbf{Округление} до требуемой разрядности.
\end{enumerate}

\paragraph*{Пример умножения}
Перемножим $A = 1.5 = 1.1_2 \times 2^0$ и $B = 2.5 = 1.01_2 \times 2^1$ в float32.

Мантиссы: $M_a = 1.1_2 = 1.5_{10}$, $M_b = 1.01_2 = 1.25_{10}$.
Произведение мантисс: $1.1_2 \times 1.01_2 = 1.111_2$? Проверим:
\[
1.1_2 \times 1.01_2 = (1 + 1/2) \times (1 + 1/4) = 1.5 \times 1.25 = 1.875 = 1.111_2
\]
(вычисления в двоичной: $1.1 \times 1.01 = 1.1 \times (1 + 0.01) = 1.1 + 0.011 = 1.111$). Результат $= 1.111_2 \times 2^{0+1} = 1.111_2 \times 2^1 = 11.11_2 = 3.75_{10}$. Нормализация не нужна ($1 \le 1.111 < 2$).

\subsubsection{Деление}
\[
X / Y = (M_x / M_y) \times 2^{E_x - E_y}
\]
\paragraph*{Алгоритм деления}
\begin{enumerate}
    \item \textbf{Деление мантисс}: $M_x / M_y$. Поскольку обе мантиссы в $[1,2)$, частное лежит в $(0.5, 2)$. Если частное $< 1$, потребуется нормализация (сдвиг влево).
    \item \textbf{Вычитание экспонент}: $E_x - E_y$.
    \item \textbf{Нормализация и округление}.
\end{enumerate}

\paragraph*{Пример деления}
Разделим $A = 3.0 = 1.1_2 \times 2^1$ на $B = 1.5 = 1.1_2 \times 2^0$ (заметим, оба имеют одинаковую мантиссу). Частное мантисс $1.1_2 / 1.1_2 = 1.0_2$. Экспонента $1 - 0 = 1$. Итог $1.0 \times 2^1 = 2.0$. Точное значение.

Другой пример: $A = 1.0 = 1.0_2 \times 2^0$, $B = 3.0 = 1.1_2 \times 2^1$. Мантиссы: $M_a=1.0$, $M_b=1.1_2=1.5_{10}$. Частное мантисс $= 1.0 / 1.1_2 \approx 0.666\ldots = 0.101010\ldots_2$. Экспонента $0 - 1 = -1$. Результат $0.101010\ldots_2 \times 2^{-1} = 1.01010\ldots_2 \times 2^{-2}$ (после нормализации). При ограниченной разрядности потребуется округление.

\subsubsection{Особенности арифметики с плавающей точкой}
\begin{itemize}
    \item \textbf{Округление:} если точное значение не представимо, оно округляется к ближайшему чётному (по умолчанию).
    \item \textbf{Поглощение:} при сложении чисел с разными порядками меньший операнд может быть сдвинут до нуля. Пример: $1.0 + 10^{-8}$ для float даст $1.0$, так как точности не хватает выразить сумму (в float $\varepsilon \sim 10^{-7}$, $10^{-8}$ меньше разницы между соседними числами около $1.0$). Для double аналогично при больших различиях.
    \item \textbf{Потеря ассоциативности:} $(a+b)+c \neq a+(b+c)$ из-за округлений. Необходимо аккуратно выбирать порядок суммирования (например, от меньших к большим для повышения точности).
    \item \textbf{Сравнение:} не рекомендуется сравнивать float на строгое равенство; используют допустимую погрешность $|a-b| < \varepsilon$.
\end{itemize}



\newpage
\section{Семинар №3. Числа с плавающей точкой. }

\subsubsection*{Задача 1 (Нормализация и представление)}
Число $23.375_{10}$ представьте в нормализованной двоичной форме, а затем в формате float IEEE 754. Укажите шестнадцатеричное представление (как 8 hex-цифр).

\textbf{Решение:}
$23_{10} = 10111_2$; $0.375 = 0.011_2$ ($0.375\times2=0.75$, $0.75\times2=1.5$, $0.5\times2=1.0 \to 011$). Число $10111.011_2$. Нормализуем: $1.0111011_2 \times 2^4$.

Экспонента истинная $= 4$; $e = 4+127=131 = 10000011_2$.

Мантисса $= 0111011$, дополняем нулями до 23 бит: $01110110000000000000000$.

Знак $= 0$.

Итого: \texttt{0 10000011 01110110000000000000000}.

Переведём в hex: \texttt{0100 0001 1011 1011 0000 0000 0000 0000} = \texttt{0x41BB0000}.

\subsubsection*{Задача 2 (Декодирование float)}
Дано 32-битное число в hex: \texttt{0xC0A00000}. Определите, какое десятичное число оно представляет (float).

\textbf{Решение:}
Hex: \texttt{C0A00000} $\to 1100\,0000\,1010\,0000\,0000\,0000\,0000\,0000_2$.

Знак $S=1$ (отрицательное).

Экспонента $e = 10000001_2 = 129$. Истинная $E = 129 - 127 = 2$.

Мантисса: дробная часть = $010\,0000\ldots0 = 0.01_2$ (в двоичном) или $1.01_2$ (с ведущей единицей). $1.01_2 = 1 + 0.25 = 1.25$.

Значение $= -1.25 \times 2^2 = -5.0$. Ответ: $-5.0$.

\subsubsection*{Задача 3 (Машинный эпсилон)}
Объясните, почему в формате float не могут быть представлены точно числа $0.1$ и $0.2$. Вычислите, какой именно $\varepsilon$ для float.

\textbf{Решение:}
$0.1_{10}$ в двоичной — периодическая дробь: $0.0001100110011\ldots_2$. При усечении до $23+1$ бита мантиссы происходит округление. Число приближается с точностью до $\varepsilon$.

Машинный $\varepsilon$ для float: $2^{-23} \approx 1.1920929 \times 10^{-7}$.

(В Python: \texttt{import sys; sys.float\_info.epsilon} даст для double; для float в C: \texttt{FLT\_EPSILON} = $1.1920929\times10^{-7}$.)

Любые числа представимы с относительной погрешностью $\le \varepsilon/2$ (при округлении к ближайшему).

\subsection*{Задача 4 (Ручное сложение float)}
Сложите два числа в формате float32 (одинарная точность). Даны десятичные числа:
\[
A = 10.0,\quad B = 0.625.
\]
Выполните сложение, имитируя процесс выравнивания порядков, сложения мантисс и нормализации. Получите результат в десятичном виде. Убедитесь, что он точен (нет округления).

\textbf{Решение:}
\begin{align*}
A &= 10.0 = 1010.0_2 = 1.010_2 \times 2^3. \quad E_a = 3,\; M_a = 1.010_2.\\
B &= 0.625 = 0.101_2 = 1.01_2 \times 2^{-1}. \quad E_b = -1,\; M_b = 1.01_2.
\end{align*}
Выравнивание до $E=3$: сдвигаем $M_b$ вправо на $3-(-1)=4$ разряда:
\[
M_b' = 0.000101_2.
\]
Сложение мантисс (с учётом скрытого бита, биты после запятой):
\[
\begin{array}{r}
  1.010000_2 \\
+ 0.000101_2 \\
\hline
  1.010101_2
\end{array}
\]
Сумма $= 1.010101_2 \times 2^3 = 1010.101_2 = 10 + 0.5 + 0.125 = 10.625_{10}$.
Точно представимо в float32. Округления не требуется.

\subsection*{Задача 5 (Поглощение)}
Даны два числа одинарной точности:
\[
A = 1.0 \times 10^8,\quad B = 1.0.
\]
Покажите, что при сложении $A + B$ результат остаётся равным $A$ из-за поглощения. Подсказка: выразите $A$ в двоичной нормализованной форме и оцените разницу порядков.

\textbf{Решение:}
$A \approx 100\,000\,000$. Переведём в двоичную: $100\,000\,000 \approx 101111101011110000100000000_2$, что примерно $1.011111\ldots_2 \times 2^{26}$. ($2^{26}=67\,108\,864$). Истинная экспонента $E_a \approx 26$ (точнее $26$, т.к. $2^{26} = 67\,108\,864$, а $10^8/67M \approx 1.49$, так что $E_a=26$).
$B = 1.0 = 1.0 \times 2^0$, $E_b = 0$.

Разница порядков: $26 - 0 = 26$. Для float32 мантисса с учётом скрытого бита имеет $24$ значащих бита. При сдвиге $B$ для выравнивания до $2^{26}$ его мантисса сдвинется вправо на $26$ разрядов. Так как $26 > 24$, все значащие биты $B$ будут вытеснены за пределы разрядной сетки. Результат сложения мантисс не изменит $A$. Следовательно, $A + B = A$, поглощение.

\subsection*{Задача 6 (Катастрофическая потеря точности)}
Даны два числа $X = 1.23456789$ и $Y = 1.23456780$. При вычитании $X - Y$ в десятичной арифметике получается $9 \times 10^{-8}$. Предположим, что эти числа хранятся в float32. Оцените относительную погрешность результата, если исходные числа представимы с точностью $\varepsilon \approx 1.2\times10^{-7}$. Сколько верных значащих цифр останется?

\textbf{Решение:}
Абсолютная погрешность каждого числа порядка $\varepsilon/2 \cdot |X| \approx 0.6 \times 10^{-7} \times 1.23 \approx 7.4 \times 10^{-8}$. Разность двух таких чисел имеет абсолютную погрешность до $2 \times 7.4\times10^{-8} \approx 1.5\times10^{-7}$, что превышает саму разность $9\times10^{-8}$. Следовательно, результат может оказаться полностью неверным (относительная погрешность $>100\%$). Значимых цифр практически не останется. Это катастрофическая потеря точности.

\subsection*{Задача 7 (Умножение с округлением)}
Перемножьте числа $A = 1.75$ и $B = 2.5$ в формате с 4 битами дробной части мантиссы (аналог float32, но усечённый). Получите точный двоичный результат, а затем округлите до 4 бит после запятой.

\textbf{Решение:}
$A = 1.75 = 1.11_2 \times 2^0$, $B = 2.5 = 1.01_2 \times 2^1$.
Произведение мантисс: $1.11_2 \times 1.01_2$. Вычислим:
$1.11_2 = 1 + \frac{3}{4}$; $1.01_2 = 1 + \frac{1}{4}$. Произведение $= 1.75 \times 2.5 = 4.375$.
В двоичной: $4.375 = 100.011_2 = 1.00011_2 \times 2^2$.
Экспоненты: $0+1=1$, но после возможной нормализации: произведение мантисс $= 1.11 \times 1.01$. Посчитаем в столбик:
\[
\begin{array}{c@{\quad}c@{\quad}c@{\quad}c@{\quad}c}
  & 1 & 1 & 1 & \quad (\text{множимое } 1.11_2)\\
\times & 1 & 0 & 1 & \quad (\text{множитель } 1.01_2)\\\hline
  & 1 & 1 & 1 & \\
1 & 1 & 1 & & \\
1 & 1 & 1 & & \\\hline
1 & 0 & 0 & 0 & 1\,1\\
\end{array}
\]
Получилось $10.0011_2$ (пять бит). Т.е. мантисса $10.0011_2$, что равно $1.00011_2 \times 2^1$ (после сдвига вправо на 1). Учитывая экспоненты, итоговая истинная экспонента: $E_a+E_b=1$, плюс 1 от нормализации = $2$. Итак, точный результат $1.00011_2 \times 2^2$.

Теперь округляем до 4 бит дробной части (после запятой в $1.f$). Имеем $1.0001\,1$ (5 бит дробной части: 00011). Округляем к ближайшему чётному. Третий бит (взвешенный $2^{-3}$) равен 1, а остальные биты отбрасываемой части = 1 (единица в $2^{-4}$). Так как отбрасываемая часть больше половины ($0.00011 > 0.00010$), округляем вверх: $1.0001 + 0.0001 = 1.0010$. Результат: $1.0010_2 \times 2^2 = 100.10_2 = 4.5_{10}$. Точное значение $4.375$ было округлено до $4.5$.

\subsection*{Задача 8 (Деление с денормализованным результатом)}
Разделите $A = 1.0 \times 2^{-126}$ (наименьшее нормализованное float32) на $2$. Объясните, что произойдёт.

\textbf{Решение:}
$A = 1.0 \times 2^{-126}$. Делим на $2 = 2^1$: экспонента уменьшается на 1, получаем $E = -127$. Но минимальная истинная экспонента для нормализованных чисел равна $-126$. Результат становится денормализованным. Мантисса $1.0$ сдвигается вправо на 1 (для нормализации в денормализованном виде скрытая единица становится 0): $0.1_2 \times 2^{-126} = 1.0 \times 2^{-127}$? Точное значение $= 2^{-127}$. В формате float32 оно представимо как денормализованное число с $e=0$, $f = 1000\ldots0_2$ (что соответствует $0.1_2 \times 2^{-126} = 2^{-127}$). Таким образом, плавный переход к нулю обеспечивается денормализованными числами.

\subsection*{Задача 9 (Особые значения и арифметика)}
Что будет результатом следующих выражений в арифметике float:
\begin{enumerate}
    \item[a)] $1.0 / 0.0$
    \item[b)] $0.0 / 0.0$
    \item[c)] $\infty - \infty$
    \item[d)] $10^{38} * 10^{38}$ (float)
\end{enumerate}
Определите битовое представление бесконечности и NaN.

\textbf{Решение:}
\begin{enumerate}
    \item[a)] $+\infty$ ($e=255$, $m=0$, $S=0$).
    \item[b)] NaN ($e=255$, $m\neq0$).
    \item[c)] NaN.
    \item[d)] Переполнение: $\infty$ (float max $\approx 3.4\times10^{38}$, $10^{38}\times10^{38} >$ max $\to \infty$).
\end{enumerate}
Битовое представление $+\infty$: \texttt{0 11111111 000\ldots0}.

NaN: \texttt{0 11111111} (любая ненулевая мантисса), обычно \texttt{0 11111111 1000000\ldots0}.

\subsection*{Задача 10 (Поглощение и неассоциативность)}
Покажите на примере чисел в формате float (одинарная точность), что $(1.0 + 10^{-8}) + 10^{-8}$ может отличаться от $1.0 + (10^{-8} + 10^{-8})$.

\textit{Указание:} учтите машинный эпсилон и выравнивание порядков.

\textbf{Решение:}
$\varepsilon_\text{float} \approx 1.19\times10^{-7}$.

$10^{-8} = 1\times10^{-8} \approx 2^{-26.6}$. При сложении $1.0$ и $10^{-8}$: для выравнивания нужно сдвинуть мантиссу $10^{-8}$ на $\approx 26$ разрядов вправо, а мантисса float хранит только 23 бита дробной части. Поэтому $10^{-8}$ выйдет за пределы разрядной сетки, и сумма будет равна $1.0$. Тогда $(1.0+10^{-8})+10^{-8} = 1.0+10^{-8} = 1.0$.

С другой стороны, $10^{-8}+10^{-8} = 2\times10^{-8}$, что всё ещё меньше $\varepsilon$ относительно $1.0$? $2\times10^{-8} \approx 2^{-25.5}$, тоже меньше $\varepsilon$, и при сложении с $1.0$ даст $1.0$. В этом примере разницы нет — он неудачен. Лучше взять $1.0 + 10^{-7}$: $10^{-7}$ немного меньше $\varepsilon$, результат может быть $1.0$. А $(1.0+10^{-7})+10^{-7} = 1.0$, $1.0+(2\times10^{-7})$ — если $2\times10^{-7} > \varepsilon/2$? $\varepsilon=1.19\times10^{-7}$, $2\times10^{-7}\approx1.68\varepsilon$. По правилам округления ближайшее к $1+2\times10^{-7}$ представимое число может быть $1+2\varepsilon$? Покажем классический пример: $a=1.0$, $b=c=10^{-16}$ для double: $\varepsilon=2\times10^{-16}$, тогда $(1+10^{-16})+10^{-16} = 1.0$, а $1+(2\times10^{-16})=1.0000000000000004$ (разница порядка ULP). Итак, неассоциативность налицо.

Итог: из-за потери точности при выравнивании порядков результаты различаются.
